\section{Related work}

\subsection{Portfolio Optimization and View Uncertainty}
Portfolio choice is difficult because the inputs are noisy, not because the optimization problem is hard \citep{markowitz1952portfolio}. The optimizer amplifies input error. Small errors in expected returns can produce large and unstable weights \citep{jobson1980estimation, best1991sensitivity}. Noise in expected returns often dominates performance out of sample, and even modest misspecification can lead to unintuitive allocations and unstable tail risk exposures \citep{merton1980estimating, chopra1993effect}. Two research lines address these issues. One line shrinks inputs through Bayes Stein mean shrinkage \citep{jorion1986bayes}, high dimensional covariance estimators \citep{ledoit2003improved, ledoit2004well}, and Bayesian priors \citep{avramov2002stock, pastor2000portfolio}. Another line regularizes the optimizer through resampling \citep{michaud1989markowitz} and robust formulations that account for model misspecification \citep{garlappi2007portfolio, bertsimas2011theory}. Many studies conclude that without strong prior structure or disciplined shrinkage, plug in mean variance optimization can underperform a simple equal weight rule \citep{demiguel2009optimal}.

Black and Litterman provide a Bayesian approach that combines an equilibrium prior with investor views \citep{black1992global}. The model uses a view uncertainty matrix to determine how strongly each view influences the posterior. This structure makes the framework well suited as an interface for external forecasts, including forecasts produced from text based models. The practical difficulty is that most applications still rely on subjective choices for both the view values and their uncertainty. Prior work explains how to express absolute and relative views \citep{he2002intuition}, and later papers clarify and extend the framework \citep{satchell2007demystification, ccela2021mean}. Industry practice often maps user reported confidence into the view uncertainty matrix through heuristic scaling rules \citep{idzorek2007step}. These choices are hard to audit, brittle across assets and horizons, and difficult to scale.

This gap matters more when views come from LLMs, since the views originate from unstructured information and the uncertainty is not directly observed. Our work follows the same goal as the prior literature on disciplined inputs, but we target the specific interface needed for optimizer level use. We elicit a forecast distribution from an LLM by repeating the same query many times under a fixed prompt, and we interpret the dispersion of the forecasts as model uncertainty about the view. We set the view to the average forecast and set view uncertainty to the forecast variance, or the forecast covariance when we elicit joint views. We treat this uncertainty as model uncertainty rather than sampling error of the average, so it does not shrink with the number of repeated queries. This mapping replaces ad hoc tuning with an auditable procedure that scales across assets, horizons, and forecasting models.

\subsection{The Missing Optimizer Link}
A parallel literature in asset pricing and corporate finance extracts forward looking information from unstructured text. Early studies show that media tone and disclosure language predict investor reactions and returns \citep{tetlock2007giving, tetlock2008more, loughran2011liability}. Surveys document the maturation of text as data, from dictionaries to machine learning and embeddings \citep{gentzkow2019text, hwang2025deep}. Large language models extend this toolkit. They can synthesize news, filings, and earnings call dialogue into forecasts at scale, and recent evidence links LLM based assessments to fundamentals and returns \citep{lopez2023can, liu2023fingpt, lee2025your}.

Two frictions keep this literature separate from formal allocation. First, most applications remain task level rather than optimizer level. Researchers often report point scores or rankings, but a portfolio optimizer needs a clear view value and a defensible measure of view uncertainty. When portfolios are formed, signal strength is usually implicit, so the optimizer cannot control exposure to uncertain views. Second, evaluation practices often fall short of financial econometric standards. General purpose LLMs are trained on corpora that include financial text, so backtests without verified knowledge cutoffs and strict time stamping risk look ahead and data contamination. This compounds classic concerns about multiple testing and data snooping \citep{bender2021dangers, drinkall2024time, kong2025fusing}.

Our work connects these threads by providing a concrete mapping from LLM outputs to the inputs required by a Bayesian optimizer. We treat the LLM forecast average as the view and the forecast dispersion as view uncertainty, and we enforce strict knowledge cutoffs and time stamped information sets to prevent contamination. This design allows language based forecasts to enter portfolio optimization on disciplined terms and enables direct evaluation of how uncertainty calibration affects allocation outcomes.

